\chapter{总结与展望}

\section{全文工作总结}

本文围绕无机晶体逆向生成中的两个推理阶段问题展开研究。第一个问题是：在 MatterGen 的固定条件引导下，如何在有限推理预算中提高目标属性生成精度，同时控制候选偏离参考轨迹带来的质量风险。第二个问题是：如何在晶体扩散生成过程中引入局部物理反馈，改善生成结构在原子力层面的代理物理一致性。两个问题分别对应条件属性控制和局部结构物理一致性，不能由单一指标或单一推理机制代替。

全文以冻结的 MatterGen 作为共同基础，不重新训练主干权重，也不改变原有三字段生成定义。研究路线先从 C0 基线出发，分别分析条件得分残差和预测干净结构的阶段可靠性；随后形成两条相互独立的推理增强路线；最后使用匹配种子、冻结代理指标、配对统计和跨代理评价进行统一分析。整体技术路线可概括为：MatterGen baseline → 条件属性多分支推理与安全回退、后期神经力场反馈 → 统一评价和证据边界审计。两条路线共享实验设计原则，但不是串联模块，也没有被验证为一个联合最优模型。

第一条路线的研究过程并非直接假定多分支有效。研究先考察固定 CFG 和多字段残差驱动的在线 Adaptive CFG，再用反事实 Oracle 判断是否存在可利用的引导空间，并通过风险、字段和时间阶段实验定位在线选择的不稳定来源。由于当前中间状态特征不能可靠判断何时应当偏离 C0，研究问题被重新表述为：保存一条可精确恢复的参考轨迹，在有限预算内展开少量候选，并使用终点属性和质量约束作出选择。由此形成参考轨迹保留的预算约束多分支条件引导与安全回退方法，并以 Fixed-K2 作为最终确认策略。

第二条路线针对属性条件满足但结构仍可能具有较高代理受力的问题。研究先在独立诊断样本上比较预测干净结构与终态 MaxF 的阶段排序一致性，再选择后期窗口调用 Matter\allowbreak Sim，将笛卡尔力映射到周期分数坐标，并通过 position predictor 的实际 score 系数注入有界位置修正。该方法不改动原子类型和晶胞的 MatterGen 更新，也不把神经势当作替代生成器；异常时保留原始 score 并执行安全回退。最终形成基于阶段可靠性校准的后期神经力场引导扩散方法 RC-NFGD。

因此，本文的主要贡献不是把所有探索路径都归纳为正结果，而是把两个不同层面的推理控制问题转化为可复现、可比较和可审计的算法接口。第3章回答条件属性控制问题，第4章回答局部代理受力反馈问题，第5章进一步区分确认结果、机制结果、负结果和证据边界。

\section{主要研究成果}

\subsection{条件属性生成方面的研究成果}

本文提出参考轨迹保留的预算约束多分支条件引导与安全回退方法（Reference-Preserved Budget-Constrained Multi-Branch Guidance with Safety-Constrained Fallback）。方法在共享前缀后保存精确 C0 参考状态和随机数状态，固定展开两个候选后缀，并在终点依据属性改善和冻结质量条件决定接受或参考回退。其核心是把难以提前预测的``何时偏离''转化为有限候选的终点验证，同时明确控制候选数量和额外 score 调用。

在 C1-128 完全独立配对种子上，Fixed-K2 将 Property MAE 从 0.034796 降至 0.026332，相对降低 24.32\%，并保持冻结的 E-hull、Stable、NUS 和 Validity 代理护栏。该结果支持以下范围内的结论：在当前 MatterGen checkpoint、目标属性、候选库和终点评价协议下，参考轨迹保留、固定双候选和安全回退可以改善代理属性误差。该方法约使用 C0 的 2.2 倍生成预算，因而属于 inference-time compute--quality trade-off，而不是推理加速。

学习型 Linear-K2 的结果必须与上述正结果同时报告。Phase B 留出测试只有区间跨零的方向性结果，C1-128 中 Linear-K2 的 Property MAE 为 0.027522，高于 Fixed-K2 的 0.026332，未达到预注册的候选分配优势门槛。因此，当前证据支持固定候选组合和终点回退，不支持``基于当前分支前特征的学习型样本级候选分配优于 Fixed-K2''。这一负结果限定了方法结论，也说明更复杂的在线控制并不会因为增加模型形式而自动获得稳定收益。

综合创新点1的证据，可以得到四点认识：MatterGen 条件生成中存在可利用的 guidance headroom；仅凭当前中间残差和短视野信号难以稳定识别有利偏离；参考保留、多分支验证和终点回退能够把风险限制在固定预算内；在当前数据和候选库下，简单 Fixed-K2 比 learned allocation 更稳健。上述结论不应写作在线 Adaptive CFG 已得到稳定正向验证，也不外推到其他属性、材料体系或生成模型。

\subsection{局部结构物理一致性方面的研究成果}

本文提出基于阶段可靠性校准的后期神经力场引导扩散方法（Reliability-Calibrated Late-stage Neural Force-Guided Diffusion，RC-NFGD）。方法先用预测干净结构的 MaxF 与最终 MaxF 排序相关性确定后期介入窗口，再通过周期坐标映射和 score 系数补偿把 Matter\allowbreak Sim 力反馈送入 position predictor，并使用有界修正和异常回退控制数值幅度。阶段诊断在当前任务中选择 (t=0.02)，其 Spearman 相关系数为 0.925587；这只是当前模型和任务的可靠性依据，不是所有扩散阶段的普遍规律。

Formal256 的冻结结果显示，RC-NFGD 使 Matter\allowbreak Sim MaxF、Mean Force 和 RMSD 分别降低 29.81\%、29.72\% 和 14.17\%，同时 Stable 保持 80.86\%，Validity 保持 100\%。P0、Formal\allowbreak 32 和 Formal256 三个分离 cohort 的 MaxF 方向一致，说明后期代理力反馈在当前实验设置下具有可复现性。CHGNet 未参与在线力引导，在同一 256 对结构上观察到 MaxF 降低 14.36\%、Mean Force 降低 9.59\%，为跨神经势的均值方向提供了额外支持。

该结果必须与不利指标一并解释。Formal256 Property MAE 从 0.009757 增至 0.009868，约恶化 1.14\%；CHGNet 的部分高分位尾部指标呈混合结果；等预算生成后 POST 在所测 Matter\allowbreak Sim 终态受力指标上更低。因此，RC-NFGD 的正式结论是：在当前冻结任务和机器学习势代理协议下，后期在线神经力反馈改善了若干代理局部结构一致性指标，并得到 CHGNet 均值方向支持。它不是 DFT 物理验证，也不是对生成后松弛的支配性结论。

创新点2还保留了两个重要的机制边界。真实力方向相对于匹配尺度的随机、打乱和反向控制具有更好的受力结果，说明方向信息具有实用价值，但消融路径仍有重放限制；去除有界修正后 Matter\allowbreak Sim 指标更低，故 bound 的定位是保守的步长和异常控制机制，而不是已被证明的性能核心。由此，RC-NFGD 应被理解为一个可审计的后期采样反馈接口，而不是完整的结构优化器。

\subsection{两项成果的关系与总体回答}

两项创新分别回答两个研究问题。问题1的回答是：在允许增加约 2.2 倍生成预算时，保存精确 C0、展开固定双候选并进行终点约束选择，可以在 C1-128 上改善代理属性误差；当前 Linear-K2 没有显示额外优势。问题2的回答是：在阶段可靠性诊断支持的后期窗口内，Matter\allowbreak Sim 力反馈可以接入 MatterGen 逆扩散，并在 Formal256 上改善代理 MaxF、Mean Force 和 RMSD；这一改善伴随轻微属性代价，且仍局限于代理评价。

联合兼容性实验的状态为 \texttt{JOINT\_SYNERGY\ =\ NOT\_SUP\allowbreak PORTED}。因此，本文将两条路线作为相互独立的推理阶段增强方法，而不是写成``创新点1之后接创新点2''的统一模型。当前研究更适合表达为：条件属性准确性和局部结构一致性是相关但不等价的目标，需要分别建模、分别设定预算和分别验证。

\section{研究局限}

第一，证据仍停留在机器学习势和属性代理层面。Matter\allowbreak Sim 同时参与 RC-NFGD 的力反馈和主要力评价，存在同源代理自洽风险；CHGNet 提供了不参与在线力引导的交叉评价，但仍然是机器学习势，且其尾部指标并非全部同向改善。全文没有进行 DFT 单点力、DFT 能量、声子、长时间分子动力学或实验合成验证，\texttt{DFT\_VERIFIED\ =\ false}。因此，代理受力下降不能被改写为真实热力学稳定性、动力学稳定性或可合成性已经得到证明。

第二，创新点1的计算代价和选择依赖需要明确。Fixed-K2 约使用 C0 的 2.2 倍 MatterGen score 调用，并依赖终点评价器、冻结候选库和接受规则。其零属性损失来自终点选择与参考回退设计，不表示每个候选分支都成功；当前实验也没有证明它优于所有等预算的通用 Best-of-N 策略。换言之，Fixed-K2 是有明确成本的质量---计算权衡，而不是无额外代价的普适改进。

第三，Linear-K2 的失败只约束当前的特征、训练规模、候选库和线性分配模型。179 个分支前特征及其线性交互没有支持稳定的跨 cohort 候选分配，但这不能推出任何学习型分配器都不可能有效。类似地，Fixed-K2 的候选身份、分支点和终点阈值来自当前 MatterGen 与当前目标属性，不能直接外推到其他 checkpoint、属性、组成范围或晶体生成模型。

第四，RC-NFGD 存在目标之间的权衡和竞争基线边界。Formal256 的 Property MAE 轻微恶化，说明低代理受力并不自动带来更好的条件属性；CHGNet 尾部混合结果说明均值改善不能代表极端样本风险全部下降；等预算 POST 在所测终态受力上更强，因而不能写成在线方法优于生成后修正。有界修正的必要性也没有被消融确立，其作用应限定为保守的数值控制。

第五，两项创新的联合效果和更广泛迁移性仍未解决。当前 Adaptive CFG 与力反馈的兼容性实验未达到双保持门槛，且没有直接确认 Fixed-K2 与 RC-NFGD 的完整笛卡尔积组合。方法还依赖当前的阶段窗口、神经势、目标属性、采样步数和结构规模；跨任务使用时必须重新进行阶段诊断、候选校准、预算设计和确认实验，不能从本文结果直接推导部署最优策略。

此外，配对种子设计提高了方法间比较的内部可比性，但不等同于对整个材料分布的无偏抽样。C1-128 和 Formal256 的确认结果分别服务于两个冻结任务，机制消融和兼容性实验的样本量更小，因而不同证据层级不能简单合并为一个总体效应。20 000 次配对 bootstrap 描述的是给定 cohort 和指标下的抽样不确定性，并不替代跨数据集、跨 checkpoint 或跨实验条件的复现。后续研究还需要报告结构规模、元素分布、属性难度和计算预算的分层结果，以判断平均改善是否由特定子群驱动，并检查回退率、尾部风险和属性代价在不同子群中是否保持同一方向。

\section{未来工作展望}

\subsection{以配对 DFT 单点力验证代理结论}

最优先的工作是对 Formal256 中冻结的 C0 与 RC-NFGD 配对结构进行随机或分层随机抽样，开展配对 DFT single-point force validation。第一阶段应优先比较 DFT 原子力的 MaxF 和 Mean Force，并沿用相同 seed 的配对统计；在计算资源允许时，再补充 single-point energy。形成能和 E-hull 还需要参考相、竞争相和一致的计算泛函，实施复杂度更高，不应在没有实际测量的情况下预先承诺成本或效果。只有建立了 DFT 层面的配对证据，才能判断代理力改善在更高精度势能面上是否保持。

\subsection{研究更可靠的候选分配与更低成本分支}

对于创新点1，后续应研究具有更强跨 cohort 泛化能力的状态表征、置信度估计和校准选择机制，而不是简单增加 Linear-K2 的复杂度。候选分支还可以探索共享计算、早期剪枝、分层预算和低成本终点评价，以降低 Fixed-K2 的 2.2 倍开销；这些方向需要在新的预注册协议下验证，不能视为当前结果的自然延伸。

\subsection{扩展属性、材料体系和生成模型}

未来可在多个属性、冲突目标和不同组成范围下测试参考保留框架，例如将属性误差、稳定性代理和结构质量写成明确的多目标门槛或 Pareto 选择问题。同时，应在其他 MatterGen checkpoint、其他晶体扩散模型及更大或更小结构规模上重新校准分支点与终点评价，以判断方法的迁移性。扩展实验的重点是检验适用范围，而不是预设当前效应量会保持不变。

\subsection{发展多物理量反馈和在线---后处理组合}

对于创新点2，可以研究能量、力和应力的多物理量反馈，但应先建立各信号在不同扩散阶段的可靠性和冲突处理规则。鉴于当前 POST 在终态代理力上更强，合理的后续问题是比较``后期在线反馈 + 轻量终态松弛''的预算组合，并明确区分生成分布保持、终态受力和属性准确性三个目标。任何组合策略都需要独立的预算、配对种子和预注册判据。

\subsection{面向实验可验证性的候选筛选}

在完成高精度计算验证后，可进一步把生成候选送入缺陷、温压条件和合成可行性筛选，再选择少量结构进行实验合成与表征。该阶段的重点应是建立从生成、代理筛选、高精度计算到实验反馈的可追溯链路，而不是把当前代理结果直接解释为实验候选已经确定。对于硕士论文后续工作，优先级仍应放在 DFT 配对验证和跨任务复现，实验路线可作为更长期的延伸。

\section{本章小结}

本文围绕 MatterGen 的条件属性控制和局部结构物理一致性，形成了两条独立的推理阶段优化路线。创新点1以参考轨迹保留、Fixed-K2、终点约束和安全回退改善 C1-128 上的代理 Property MAE；学习型 Linear-K2 的额外优势未得到支持。创新点2以阶段可靠性校准、后期 Matter\allowbreak Sim 力反馈和周期坐标映射改善 Formal256 上的代理 MaxF、Mean Force 和 RMSD；CHGNet 对均值力指标给出方向一致支持，但 Property MAE 存在轻微代价。

全文同时保留了关键边界：两项创新没有得到已验证的联合协同，等预算 POST 在所测终态受力上更强，有界修正不是已确认的性能核心，且没有 DFT 或实验验证。由此，本文形成的最终结论是：Fixed-K2 在冻结代理属性协议下得到支持，RC-NFGD 在冻结代理局部结构协议下得到支持；更强的真实物理、普适迁移和联合最优命题仍需新的预注册研究。后续工作的首要任务是配对 DFT 单点力验证，其次是降低分支成本、改进候选分配并扩大跨任务评估。
